\section{Background and Related Work}
\label{sec:rw}

Our work builds on four fundamentals: (i) developer turnover and disengagement in open source, (ii) burnout and the Job Demands-Resources model, (iii) affect mined from software repositories, and (iv) psycholinguistic markers of identification and recognition.
\begin{description}[leftmargin=0.3cm]
\item[Turnover and Disengagement in Open Source]
Open-source communities are organized around a small core of developers surrounded by a large periphery~\cite{joblin2017classifying,joblin2015developer}, and their survival depends on a few key contributors, as captured by the truck factor~\cite{avelino2016novel}. Avelino \etal~\cite{avelino2019abandonment} showed that the abandonment of truck-factor developers is common and threatens project survival. Lin \etal~\cite{lin2017developer} applied survival analysis to developer turnover, linking it to activity patterns, while Iaffaldano \etal~\cite{iaffaldano2019why} and Calefato \etal~\cite{calefato2022will} studied the breaks of core developers and showed that many of them eventually come back. Miller \etal~\cite{miller2019people} found that contributors often disengage for reasons external to the project, such as changing jobs, and motivations shift over time~\cite{gerosa2021shifting}. These studies explain turnover through what developers \textit{do}, \ie activity, tenure, and position, but not through how they \textit{feel} and relate to their community before leaving. Moreover, they rarely distinguish developers who leave a project from those who leave open source.

\item[Burnout and the Job Demands-Resources Model]
Burnout is a syndrome of exhaustion, cynicism, and reduced efficacy caused by unmanaged chronic workplace stress~\cite{maslach2016burnout}. Tulili \etal~\cite{tulili2023burnout} mapped 92 software engineering studies on burnout, most of which aim at early prediction through sentiment analysis. The dominant lens is the JD-R model~\cite{demerouti2001job,bakker2007job}, which links demands to exhaustion through a \textit{health-impairment} process and resources, such as recognition and social support, to engagement through a \textit{motivational} process. Large-scale surveys linked burnout to belonging and organizational culture~\cite{trinkenreich2023model} and to actual attrition~\cite{trinkenreich2024attrition}, while Raman \etal~\cite{raman2020stress} pointed to unhealthy interactions as a source of stress in open source. These studies rely on self-reports collected at one point in time, which cannot follow developers up to the moment they leave.

\item[Affect in Software Repositories]
Research on developers' affect~\cite{novielli2019sentiment,sanchez2019taking} has mined sentiment and emotions from issues, commits, and Q\&A sites~\cite{murgia2014developers,calefato2018sentiment,imran2024uncovering}, relating them to productivity~\cite{graziotin2014happy} and issue resolution~\cite{ortu2015bullies}, and has been proposed as an alternative to work-health surveys~\cite{mantyla2016mining}. A recurring finding is that affect classifiers do not transfer across platforms: tools trained on Stack Overflow lose accuracy on GitHub~\cite{novielli2020can}, and emotion classifiers such as EmoTxt~\cite{calefato2017emotxt} reach a micro-F1 of 0.42 on GitHub comments, improved by data augmentation~\cite{imran2022data}. Most studies analyze affect at the level of messages; few follow the affect of the same developers over time, and none relates it to their disengagement.

\item[Identification and Recognition in Language]
Psycholinguistics shows that function words reveal psychological states: the use of first-person plural pronouns signals identification with a group, whereas the use of singular ones signals self-focus and, in some contexts, distress~\cite{tausczik2010psychological}. Similarly, gratitude and positive feedback express recognition, a job resource in the JD-R model~\cite{bakker2007job}, which on GitHub takes the form of thanks in replies and of reactions such as~\faThumbsUp. Despite their low cost and transparency, these markers have not been used to study disengagement in software engineering.
\end{description}

\begin{rqsummary}{\faList\ Summary of the Background and Research Gap.}
Turnover research explains who leaves through activity and structural position, burnout research relies on cross-sectional self-reports, and affect mining analyzes messages rather than developers over time. Yet, no study has investigated whether emotions, identification with the community, and received recognition change before developers disengage, nor whether they distinguish leaving a project from leaving open source.
\end{rqsummary}
